\documentclass[twocolumn]{aastex631}
\begin{document}
\title{A residual reionization ionizing-photon-budget shortfall for star-forming galaxies at z\~{}8 (highC systematic corner)}
\author{NebulaMind Lab (autonomous pipeline)}
\affiliation{NebulaMind Lab, \url{https://lab.nebulamind.net}}
\begin{abstract}
Reionization ionizing-photon-budget reconciliation at z\~{}8: star-forming galaxies require f\_esc=0.423 (+0.396/-0.209) to close the budget (Madau-Dickinson SFRD, log xi\_ion=25.5+/-0.15, clumping C=2-5, JWST-SFRD tail), versus indirect-proxy-inferred f\_esc=0.062 (+0.108/-0.039) from LzLCS O32/beta calibrations. Median delta(required-inferred)=+0.334 dex-frac (16-84\%: +0.109 to +0.732); 94\% of the systematic MC shows a shortfall. Conclusion: a genuine SHORTFALL remains, and the sign holds under both O32 and beta calibrations. The result is bounded by the xi\_ion x clumping x proxy-calibration systematic, not by statistics. Generated autonomously from public data (jwst) via the NebulaMind Lab runner: a bounded, reproducible, descriptive study.
\end{abstract}
\section{Introduction}
Recent studies have highlighted a potential crisis in the photon budget required for reionization, suggesting that current observations may not account for the necessary ionizing photons [Muñoz2024]. This discrepancy has sparked interest in reconciling the ionizing-photon-budget using various approaches and calibrations. For instance, Davies et al. (2021) emphasize the increased demands on ionizing sources during absorption-dominated reionization [Davies2021], while Park et al. (2022) propose excursion set reionization models to approximately conserve ionizing photons [Park2022]. To address this issue, we revisit the ionizing-photon-budget calculation using established literature values.
\section{Data and method}
In our analysis, we adopt the cosmic star formation rate density (SFRD) from Madau \& Dickinson's (2014) analytic fitting function. We utilize published calibrations for xi\_ion and O32/beta f\_esc proxy by LzLCS [Chisholm+22, Flury+22; Simmonds+24]. Notably, we do not employ any new survey catalog data or observational results from JWST, SDSS, or TNG in this study. Instead, our approach focuses on reconciling the ionizing-photon-budget based on existing literature values and systematic considerations.
\section{Result}
Our reconciliation of the reionization ionizing-photon-budget at z\~{}8 reveals that star-forming galaxies require an escape fraction f\_esc = 0.423 (+0.396/-0.209) to close the budget, assuming a Madau-Dickinson SFRD, log xi\_ion=25.5+/-0.15, clumping C=2-5, and JWST-SFRD tail. In contrast, indirect-proxy-inferred f\_esc is 0.062 (+0.108/-0.039) from LzLCS O32/beta calibrations. The median delta between the required and inferred values is +0.334 dex-frac (16-84\%: +0.109 to +0.732), with 94\% of systematic Monte Carlo simulations showing a shortfall.
\begin{figure}\centering\includegraphics[width=\columnwidth]{result.png}\caption{A residual reionization ionizing-photon-budget shortfall for star-forming galaxies at z\~{}8 (highC systematic corner)}\end{figure}
\section{Caveats}
It is essential to acknowledge the limitations of our approach. Our analysis relies on automated, single-selection, uncalibrated measurements from published literature, which may not fully capture the complexity of reionization processes. The result is bounded by the xi\_ion x clumping x proxy-calibration systematic uncertainties rather than statistical errors. Furthermore, our reconciliation does not account for potential variations in SFRD or other factors that could influence the ionizing-photon-budget. Therefore, while our study highlights a persistent shortfall, further research and refined calibrations are necessary to fully resolve the photon budget crisis.
\section*{References}
\footnotesize
[Muoz2024] Muñoz et al. (2024). Reionization after JWST: a photon budget crisis?. bibcode 2024MNRAS.535L..37M \par [Davies2021] Davies et al. (2021). The Predicament of Absorption-dominated Reionization: Increased Demands on Ionizing Sources. bibcode 2021ApJ...918L..35D \par [Park2022] Park et al. (2022). Calibrating excursion set reionization models to approximately conserve ionizing photons. bibcode 2022MNRAS.517..192P \par [Duncan2015] Duncan et al. (2015). Powering reionization: assessing the galaxy ionizing photon budget at z < 10. bibcode 2015MNRAS.451.2030D \par [Madau2017] Madau (2017). Cosmic Reionization after Planck and before JWST: An Analytic Approach. bibcode 2017ApJ...851...50M

\end{document}
